% Part: sets-functions-relations
% Chapter: functions
% Section: composition

\documentclass[../../../include/open-logic-section]{subfiles}

\begin{document}

\olfileid[fa-IR]{sfr}{fun}{cmp}
\olsection{ترکیب توابع}

\begin{explain}
\oliflabeldef{sfr:fun:inv:sec}{در \olref[inv]{sec} دیدیم که
وارونِ~$f^{-1}$ِ !!a{bijection}~$f$ خود تابع است. عمل دیگری
که روی توابع انجام می‌دهیم، ترکیب است. اکنون می‌توانیم}{می‌توانیم} تابعی جدید را با
ترکیب دو تابع $f$ و~$g$ تعریف کنیم؛ یعنی نخست
$f$ و سپس~$g$ را اعمال کنیم. البته این کار تنها هنگامی ممکن است که
بردها و دامنه‌ها سازگار باشند؛ یعنی بردِ~$f$ باید زیرمجموعه‌ای از
دامنهٔ~$g$ باشد. \oliflabeldef{sfr:rel:ops:sec}{این عمل روی
توابع همانند عمل حاصل‌ضرب نسبی روی
روابط است که در \olref[rel][ops]{sec} معرفی شد.}{}

یک نمودار می‌تواند به توضیح ایدهٔ ترکیب کمک کند. در
\olref{fig:composition} دو تابع $f \colon A \to B$
و $g \colon B \to C$ و ترکیبشان~$(\comp{f}{g})$ را نشان می‌دهیم. تابع
$(\comp{f}{g}) \colon A \to C$ هر !!{element} از~$A$
را با !!a{element}ی از~$C$ جفت می‌کند. اینکه کدام !!{element} از~$C$
با !!a{element}ی از $A$ جفت می‌شود، چنین مشخص می‌گردد: با داشتن ورودی $x \in
A$، نخست تابع $f$ را بر~$x$ اعمال می‌کنیم که مقداری مانند $f(x)
= y \in B$ خروجی می‌دهد؛ سپس تابع $g$ را بر~$y$ اعمال می‌کنیم که مقداری مانند
$g(f(x)) = g(y) = z \in C$ خروجی می‌دهد.
\begin{figure}
  \olasset[2\olphotowidth]{assets/diagrams/composition.tikz}
  \caption{ترکیبِ $g \circ f$ از دو تابع $f$ و~$g$.}
  \ollabel{fig:composition}
\end{figure}
\end{explain}

\begin{defn}[ترکیب]
فرض کنید $f\colon A \to B$ و $g\colon B \to C$ تابع باشند.
\emph{ترکیبِ} $f$ با~$g$ برابر است با $\comp{f}{g} \colon A \to C$،
که در آن $(\comp{f}{g})(x) = g(f(x))$.
\end{defn}

\begin{ex}
توابع $f(x) = x + 1$ و $g(x) = 2x$ را در نظر بگیرید. چون
$(\comp{f}{g})(x) = g(f(x))$، برای هر ورودی~$x$ باید نخست
پسین آن را بگیرید و سپس حاصل را در دو ضرب کنید. پس ترکیب آن‌ها
به‌وسیلهٔ $(\comp{f}{g})(x) = 2(x+1)$ داده می‌شود.
\end{ex}

\begin{prob}
نشان دهید اگر $f \colon A \to B$ و $g \colon B \to C$ هر دو
!!{injective} باشند، آنگاه $\comp{f}{g}\colon A \to C$ نیز !!{injective} است.
\end{prob}

\begin{prob}
نشان دهید اگر $f \colon A \to B$ و $g \colon B \to C$ هر دو
!!{surjective} باشند، آنگاه $\comp{f}{g}\colon A \to C$ نیز !!{surjective} است.
\end{prob}

\begin{prob}
فرض کنید $f \colon A \to B$ و $g \colon B \to C$. نشان دهید گراف
$\comp{f}{g}$ برابر است با $R_f \mid R_g$.
\end{prob}

\end{document}
